\documentclass[12pt]{article}
\usepackage[spanish]{babel}
\usepackage[utf8]{inputenc}
\usepackage[T1]{fontenc}
\usepackage{graphicx}
\usepackage{cite}
\usepackage{float}
\usepackage{amsfonts}
\usepackage{amsmath}
\usepackage{geometry}
\graphicspath{{IMAGES}}
\newtheorem{tab}{Tabla}
\renewcommand\tablename{Tabla}
\setlength{\columnsep}{7mm}
\renewcommand\tablename{Tabla}
\renewcommand\contentsname{\'Indice}
\setlength{\topmargin}{-0.75in} \setlength{\textheight}{9.25in}
\setlength{\oddsidemargin}{0.0in} \setlength{\evensidemargin}{0.0in}
\setlength{\textwidth}{6.5in}


\begin{document}

\title{PRÁCTICA: CALOR ESPECÍFICO}
\author{Toriz R. Ángel David 201628195}
\date{10 de mayo de 2019}

\twocolumn[
\begin{@twocolumnfalse}
\maketitle
\vspace*{-0.8cm}
\begin{center}\rule{0.9\textwidth}{0.1mm} \end{center}
\begin{abstract}
\normalsize Se presentan los valores obtenidos experimentalmente del calor específico de varios materiales, con una alta coincidencia con los valores esperados.
\begin{center}\rule{0.9\textwidth}{0.1mm} \end{center}
\vspace*{0.3cm}
\end{abstract}
\end{@twocolumnfalse}
]

\section{Introducci\'on}
La energía interna de un sistema puede cambiarse por medio de trabajo o de calor, aunque de manera práctica es más fácil producir calor a partir de combustión o de electricidad que producir trabajo a partir de resortes o pesas. En consecuencia, cuando se realizaron experimentos para medir la capacidad de una sustancia para almacenar energía interna, se usó calor en vez de trabaji, y los resultados pasaron a ser llamados \emph{capacidad calorífica} de la muestra. 
Cuando un sistema absorbe calor, puede llevarse a cabo o no un cambio de temperatura, dependiendo del estado del sistema; por ejemplo, un material  su temperatura de fusión no experimenta un cambio de temperatura cuando es calentado, pero un material a una temperatura menor a la de fusión aumenta de temperatura. Si un sistema experimenta un cambio de temperatura de $T_{i}$ a $T_{f}$ durante la transferencia de Q unidades de calor, la capacidad calorífica promedio del sistema está definida como:
\begin{center}
\large{$C_{avg} = \frac{Q}{T_{f}-T_{i}}$}
\end{center}
Conforme Q y $(T_{f}-T_{i})$ se vuelven más pequeñas, esta razón se aproxima a un valor límite, conocida como la \emph{capacidad calorífica, C:}
\begin{center}
\large{$C=\lim_{T_{f}\rightarrow T_{i}} \frac{Q}{T_{f}-T_{i}}=\frac{\delta Q}{dT}$}
\end{center}
La cual es proporcional a la masa de la sustancia:
\begin{center}
\large{$c=\frac{C}{m}$}
\end{center}
ésta es la capacidad calorífica específica.\cite{ce:1}
 La cantidad $\delta Q$ de calor absorbido depende de las condiciones que hay mientras aumenta la temperatura, y se hace una diferenciación entre la capacidad calorífica a volumen V constante, $C_{v}$, y la capacidad calorífica a presión constante, $C_{p}$.
Por la Primera Ley:
\begin{center}
\large{$\delta Q= dU + pdV$}
\end{center}
$C_{p}>C_{v}$. En el caso de los sólidos, el cambio de volumen es despreciable, tal que $C_{p}=C_{v}$. Podemos entonces hallar: 
\begin{center}
\large{$C_{v}=\left(\frac{\partial U}{\partial T}\right)_{V}$}
\end{center}

De acuerdo a la teoría de Debye, la cual considera vibraciones de hasta una frecuencia $\nu_{D}$, la capacidad calorífica es:
\begin{center}
\large{$3 Nk\cdot D(T/\Theta)$}
\end{center}
Donde k es la constante de Boltzmann, N el número de átomos en el volumen considerado y $D(T/\Theta)$ es la función de Debye. Para valores grandes de $T/\Theta$, obtenemos la ley de Dulong-Petit: $C_{v}=3Nk$.
\section{Método experimental}
\subsection{M\'etodo Experimental}
Para calcular la capacidad del calorímetro, se utiliza:
\begin{center}
\large{$C_{k}=c_{w}\cdot m_{w} \frac{\vartheta_{w}-\vartheta_{m}}{\vartheta_{m}-\vartheta_{k}}$}
\end{center}
donde $\vartheta_{k}$ es la temperatura del calorímetro antes del experimento, $\vartheta_{m}$ la temperatura de la mezcla, $\vartheta_{w}$ la termperatura del agua caliente, $m_{w}$ la masa del agua y $c_{w}=4.187 \frac{J}{gK}$ es la capacidad calorífica específica del agua.

Ya que se conoce este valor, se calientan las muestras de metal en el agua hirviendo; el punto de ebullición, $\vartheta_{2}$ se determina como función de la presión atmosférica:
\begin{center}
\large{$\vartheta_{2}=100+0.0276(p-1013hPa)-0.000017(p-1013hPa)$}
\end{center}
donde p está medida en hPa y $\vartheta_{2}$ en grados celsius.

El calor específico del material se calcula con:
\begin{center}
\large{$c=\frac{(c_{w}\cdot m_{w}+C_{k})(\vartheta_{m}-\vartheta_{1})}{m_{p}(\vartheta_{2}-\vartheta_{m})}$}
\end{center}\cite{ce:2}
\subsection{Arreglo Experimental}
Se utilizaron los siguientes materiales:
\begin{itemize}
\item Calorímetro
\item Barra de aluminio
\item Masa de plomo
\item Barra de acero
\item Pocillo
\item Parrilla
\item Agua
\item Termómetro con soporte
\item Balanza
\end{itemize}

\subsubsection{Realización del experimento}
Primero debe medirse la capacidad del calorímetro. Para esto, se mide primero la masa del recipiente, $M_{p}$, con la balanza, mientras en el pocillo se calientan 100 ml de agua a alrededor de 50 grados Celsius. El recipiente se tapa y se inserta un termómetro en la abertura; ya que el agua esté a la temperatura adecuada, se vierte en el recipiente, y después de cinco minutos, se anota la lectura del termómetro; todo este proceso debe repetirse varias veces para mayor precisión.

Para los metales, el procedimiento es similar: se sumergen en agua a temperaturas elevadas (las cuales se registran), y se introducen luego en el calorímetro con el termómetro ya insertado; se dejan pasar cinco minutos y se registra la temperatura. Con la fórmula obtenida anteriormente se halla el calor específico del material.

Para la comprobación de la ley de Dulong-Petit, se hace el mismo proceso que para las barras de metal, usando ahora barras de un mol de cada metal. El calor específico obtenido se multiplica por la masa molar de cada material, y el resultado debe coincidir con la ley.
\newpage
\onecolumn
\subsection{Resultados}
Para el calor específico del calorímetro se encontró lo siguiente:
\begin{figure}[!h]
    \centering
    \includegraphics[scale=0.8]{/CE1.png}
\end{figure}

Consideraremos, pues, el calor específico del calorímetro como 68.72.

Para los metales, se obtuvieron los siguientes datos:
\begin{figure}[!h]
    \centering
    \includegraphics[scale=0.8]{/CE2.png}
\end{figure}

Con los siguientes errores relativos:
\begin{figure}[!h]
    \centering
    \includegraphics[scale=0.8]{/CE3.png}
\end{figure}

Para la parte de la Ley de Petit-Dulong, se hallaron los siguientes resultados:
\begin{figure}[!h]
    \centering
    \includegraphics[scale=0.8]{/CE4.png}
\end{figure}

Y los respectivos errores relativos:
\begin{figure}[!h]
    \centering
    \includegraphics[scale=0.8]{/CE5.png}
\end{figure}
\newpage
\twocolumn
\section{Conclusiones}
Puede observarse que en la primera parte del experimento, los valores experimentales coinciden en su mayor parte con los teóricos, siendo el mayor error relativo de poco más de 12\%. En el caso de la Ley de Petit-Dulong, ninguno de los errores relativos fue mayor de 11\%, por lo que también se obtuvieron resultados positivos en esta parte.

\bibliography{ce} 
\bibliographystyle{ieeetr}
\end{document}